%% npj Digital Medicine — Collection: Synthetic Clinical Data and
%% Privacy-Preserving Frameworks for Trustworthy Health AI
%% https://www.nature.com/collections/gggjgihcjj   (deadline 2027-07-01)
%%
%% Springer Nature LaTeX template v3.1 (December 2024), official download.
%% sn-nature is THE option for Nature Portfolio journals; do not change it.
%%
%% Build:  pdflatex main && bibtex main && pdflatex main && pdflatex main
%% NOTE: never compiled locally — this machine has no TeX toolchain. Compile on Overleaf.

\documentclass[pdflatex,sn-nature]{sn-jnl}

\usepackage{graphicx}
\usepackage{multirow}
\usepackage{amsmath,amssymb,amsfonts}
\usepackage{booktabs}
\usepackage{xcolor}

\theoremstyle{thmstyleone}
\newtheorem{theorem}{Theorem}

\raggedbottom

\begin{document}

%%============================================================%%
%% TITLE — declarative, clinical, mirrors the reference's form:
%% instrument + active verb + finding + population
%%============================================================%%
\title[Membership risk in synthetic CGM]{Membership inference identifies
concentrated re-identification risk in synthetic continuous glucose
monitoring data}

\author*[1]{\fnm{Ling} \sur{Cheng}}\email{lcheng2@ic.ac.uk}

%% TODO co-authors and affiliations
\affil*[1]{\orgdiv{TODO}, \orgname{Imperial College London},
\orgaddress{\city{London}, \country{United Kingdom}}}

%%============================================================%%
%% ABSTRACT — one paragraph, unstructured, no subheadings.
%% Target 150--220 words (two npj articles measured at 138 and 220).
%% No citations, no equations.
%%============================================================%%
\abstract{TODO --- draft last, once the Results are final. Beats, in order:
(i)~synthetic CGM is being shared and its public deposition is now required;
(ii)~we built a paired membership-inference design over a 506-subject cohort and
applied it to seven generators across two data compositions;
(iii)~risk is not uniform --- the strongest published generator exposes 24 of 26
patients above AUC 0.55 against that release's own permutation floor of 7;
(iv)~the signal is carried by absolute glucose level and by temporal ordering;
(v)~a generator whose bottleneck cannot store absolute level sits at the
permutation floor in both compositions while achieving the best fidelity of any
learned model tested; (vi)~level can be stripped and re-offset at release,
preserving shape-based glycaemic metrics.}

\keywords{synthetic data, membership inference, continuous glucose monitoring,
privacy, generative models, type~1 diabetes}

\maketitle

%%============================================================%%
%% INTRODUCTION
%% The heading is dropped at typesetting in Nature Portfolio journals --- the
%% published article opens straight into the field. Keep the \section here for
%% the submission file; production removes it.
%% ~600 words, four paragraphs, NO roadmap paragraph, NO related-work section.
%%============================================================%%
\section{Introduction}\label{sec:intro}

TODO~1. Continuous glucose monitoring is among the most valuable and least
shareable physiological records; synthetic data is the proposed route; this
collection, and a growing number of funders, \emph{require} public deposition of
the synthetic dataset itself.

TODO~2. The privacy of synthetic health data is usually argued from
nearest-neighbour heuristics, or from the absence of any copied record.
Membership inference against \emph{generative} models for physiological time
series is barely measured, and CGM is widely assumed benign.

TODO~3. Aggregate privacy metrics average away the people who matter. Rare
physiology is simultaneously what sharing is meant to preserve and what is most
identifying --- the sparse-population fairness question this collection asks.

TODO~4. What we did, what we found, and what a data custodian should take from
it. Three to four sentences. No ``the remainder of this paper''.

%%============================================================%%
%% RESULTS  — subheadings are DECLARATIVE SENTENCES stating a finding.
%% Results : Discussion : Methods should end up near 4 : 1 : 2.5.
%%============================================================%%
\section{Results}\label{sec:results}

\subsection{Membership signal in synthetic CGM is bounded between a memorisation
ceiling and a permutation floor}\label{subsec:instrument}

TODO --- Fig.~\ref{fig:design}. Establish that the measurement resolves, and
correct how such numbers are read: ``patients above AUC 0.55'' has a floor of
7--12 of 26, not ${\sim}1$, and the floor is computed per released dataset.

\subsection{Membership is recoverable for nearly every patient from glucose
alone, and most confidently for those with atypical
profiles}\label{subsec:concentration}

We applied the paired design to seven generators across the two data
compositions (Table~\ref{tab:benchmark}). Read as a raw count, the number of
patients whose membership is recoverable above \mbox{AUC 0.55} suggests that
every generator leaks: the counts range from 5 to 24 of 26. The permutation
floor removes that impression. Once each released dataset is compared against
its own floor, the generators separate sharply, and the separation does not
track generation fidelity. The strongest published generator in our comparison,
DiM-TS, exposes \textbf{24 of 26} patients from glucose alone against a floor of
\textbf{7}. DiffWave, whose fidelity is an order of magnitude worse, sits at an
excess of four patients over its floor; Diffusion-TS at $-1$; FourierDiffusion
at two. The two best-fitting learned generators in the table occupy opposite
extremes of the exposure column.

Exposure at this level is not confined to unusual patients. Of the 24 patients
DiM-TS exposes from glucose alone, \textbf{13 of 13 are matched controls} and 11
of 13 are the atypical targets --- that is, every typical patient in the design
is recoverable, and marginally more reliably than the atypical ones. A study
that reported only the count would conclude that atypical physiology confers no
additional risk.

The severity of the exposure tells the opposite story. Within the same release
the median per-patient AUC is \textbf{0.80} in the atypical arm against
\textbf{0.635} in the matched control arm, where the same quantity computed on
permuted membership sits at 0.51 and 0.50 respectively. The most exposed
patient, an atypical one, reaches \mbox{AUC 0.96}; the second, also atypical,
0.89. Atypical patients are therefore not exposed \emph{more often} under this
composition --- they are exposed \emph{more confidently}, and the threshold count
is the wrong instrument for seeing it.

TODO --- one closing sentence carrying the clinical reading: from a
glucose-only release the question is not whether a patient can be recovered but
how securely, and the patients recovered most securely are those whose profiles
least resemble the cohort's.

\begin{table}[t]
\caption{Generation fidelity and membership exposure for seven generators on two
data compositions. Context-FID is lower-is-better. ``Arm AUC'' contrasts the
atypical arm against the matched control arm; 0.5 is no separation. ``Exposed''
counts patients above \mbox{AUC 0.55} out of 26, given against that release's own
permutation floor. The floor, not zero, is the reference.}\label{tab:benchmark}
\begin{tabular}{@{}lcccccc@{}}
\toprule
& \multicolumn{3}{c}{Glucose only} & \multicolumn{3}{c}{Glucose $+$ basal insulin} \\
\cmidrule(lr){2-4}\cmidrule(lr){5-7}
Generator & Context-FID & Arm AUC & Exposed & Context-FID & Arm AUC & Exposed \\
\midrule
Copy--paste (ceiling) & 0.034 & 0.562 & --- & 0.043 & 0.680 & --- \\
\textbf{This work} & \textbf{0.038} & 0.698 & \textbf{10 vs 10} & \textbf{0.055} & 0.479 & \textbf{5 vs 5} \\
Backbone, unmodified & --- & 0.515 & 6 vs 5 & 0.055 & 0.598 & 7 vs 6 \\
DiM-TS & 0.095 & 0.562 & \textbf{24 vs 7} & 0.159 & 0.698 & \textbf{15 vs 8} \\
FourierDiffusion & 0.100 & 0.710 & 7 vs 5 & 0.194 & 0.485 & 6 vs 3 \\
DiffWave & 0.368 & 0.639 & 12 vs 8 & 0.491 & 0.740 & 13 vs 11 \\
Diffusion-TS & 0.422 & 0.828 & 7 vs 8 & 2.058 & 0.604 & 8 vs 4 \\
\botrule
\end{tabular}
\footnotetext{The copy--paste ceiling has no permutation floor computed; it is a
construction rather than a fitted model. The unmodified backbone was not scored
for fidelity on the glucose-only composition. TODO: close both gaps.}
\end{table}

\subsection{The membership signal is carried by absolute glucose level and by
temporal ordering}\label{subsec:mechanism}

TODO --- Fig.~\ref{fig:ablation}.

\subsection{Adding basal insulin protects typical patients but leaves atypical
patients exposed}\label{subsec:channel}

The second data composition separates the two arms that glucose alone had left
indistinguishable by count. Releasing basal insulin alongside glucose reduces
the number of exposed matched controls from \textbf{13 of 13 to 5 of 13}, while
the atypical arm falls only from 11 of 13 to \textbf{10 of 13}
(Fig.~\ref{fig:channel}). The protection is real but it is not distributed:
it accrues almost entirely to the patients who were already typical.

This is the opposite of what the real recordings predict. On real windows, a
single day of basal insulin alone identifies its subject \textbf{29.3\%} of the
time against a chance rate of 0.11\%, where a day of glucose alone manages
\mbox{1.5--2.6\%}. Basal insulin is a programmed schedule --- person-specific,
stable, and close in character to the absolute level that
Section~\ref{subsec:mechanism} identifies as the carrier of the membership
signal. By that logic, adding it should raise exposure, not lower it.

The resolution is on the fidelity axis, and it is why the two axes must be read
together. Every generator in Table~\ref{tab:benchmark} fits the two-channel
composition worse than the one-channel composition. The released set is
correspondingly noisier, and membership signal is lost with the fidelity. What
survives that loss is the signal that was strongest to begin with, which is
carried by the atypical patients. Adding a channel that is intrinsically more
identifying therefore lowers measured exposure --- but it does so by degrading
the release, and it leaves the most vulnerable patients where they were.

TODO --- Fig.~\ref{fig:channel} must show all three quantities on one panel:
real-data identifiability by channel, measured exposure by arm, and fidelity.
Any two of them alone support the wrong conclusion.

\subsection{Removing the model capacity that stores absolute level returns
membership signal to the permutation floor}\label{subsec:modules}

TODO --- Fig.~\ref{fig:modules}. The method appears here, as a measured result.
The recipe belongs in Methods.

\subsection{Stripping and re-offsetting absolute level preserves shape-based
glycaemic metrics}\label{subsec:mitigation}

TODO --- Table~\ref{tab:utility}.

%%============================================================%%
%% DISCUSSION — continuous prose, NO subheadings, ~800 words.
%% This is where comparison with prior work lives.
%%============================================================%%
\section{Discussion}\label{sec:discussion}

TODO. In order: what the numbers license; comparison with prior work; the
disagreement between the arm-level and per-patient readings on \texttt{d1\_c1},
on a single replicate; whether the destroyed signal is the clinically useful
signal; limitations; what a custodian should do.

%%============================================================%%
%% METHODS — LAST. Noun-phrase subheadings.
%% "Sex as biological variable" is MANDATORY Nature Portfolio reporting.
%%============================================================%%
\section{Methods}\label{sec:methods}

\subsection{Sex as biological variable}\label{subsec:sex}
TODO --- report the cohort's sex composition and state explicitly whether the
analyses are sex-stratified. Mandatory for Nature Portfolio journals.

\subsection{Cohort and data acquisition}\label{subsec:cohort}
TODO.

\subsection{Window construction and subject key}\label{subsec:windows}
TODO --- the subject key is the pair (source study, participant id); the
participant id alone is unique only within a source study.

\subsection{Definition of atypical glycaemic profiles}\label{subsec:outliers}
TODO.

\subsection{Control matching}\label{subsec:matching}
TODO.

\subsection{Paired membership-inference design}\label{subsec:design}
TODO.

\subsection{Generative models and training budgets}\label{subsec:generators}
TODO.

\subsection{Membership statistic}\label{subsec:statistic}
TODO --- including the matching of released-set size between the two arms.

\subsection{Permutation floor}\label{subsec:floor}
TODO.

\subsection{Generation fidelity metrics}\label{subsec:fidelity}
TODO.

\subsection{Glycaemic utility metrics}\label{subsec:utility}
TODO --- per-chunk distributions with interquartile ranges and a distribution
distance, not a single mean delta.

\subsection{Statistical analysis}\label{subsec:stats}
TODO.

%%============================================================%%
%% BACK MATTER — order follows the published reference.
%%============================================================%%
\backmatter

\bmhead{Data availability}
TODO --- must name the deposited synthetic dataset, its repository, its
persistent identifier and its licence. The collection requires this at
submission.

\bmhead{Code availability}
The code used to generate the results presented in this study is publicly
available at \url{https://github.com/Cranooooooo/Metabonet_CGM_MIA}.

\bmhead{Acknowledgements}
TODO.

\bmhead{Author contributions}
TODO.

\bmhead{Funding}
TODO.

\bmhead{Competing interests}
The authors declare no competing interests.

\bmhead{Supplementary information}
TODO.

\bibliography{refs}

\end{document}
